% Generated from the matching Typst scaffold by
% scripts/migrate_typst_report_to_latex.py.

% ── 4.4  Retrieval as a Fixed Invariant ──────────────────────────────────
% PURPOSE: Establish retrieval as THE central experimental control — frozen so that
%          any between-policy difference is attributable to storage, not retrieval.
% MOVE: Counter-intuitive lead emphasis — the most load-bearing design decision.
% BEATS: see ordered bullets below.
% FIGURES/TABLES: fig:architecture (cosine, fixed path).
% CITATIONS: liu2024.
% ─────────────────────────────────────────────────────────────────────────

\section{Retrieval as a Fixed Invariant}

\begin{todobox}
Expand the beats below into prose. This is the linchpin of internal validity: state plainly that "retrieval held constant" is what licenses the causal reading.
\end{todobox}

% - Retrieval scoring = pure cosine similarity, IDENTICAL across all six conditions
%   (Invariant #5) — no bonuses, no penalties, no policy-specific reranking; enforced in
%   `retriever.py::shared_retrieve`.
% - top_k = 5 items returned per query; same-repo retrieval only (Invariant #16),
%   so a sequence never pulls memories from a foreign repository.
% - Injection order is relevance-sorted with the BEST item placed LAST (Invariant #6) —
%   the Lost-in-the-Middle fix, motivated by liu2024's finding that models attend most to
%   the ends of long contexts.
% - Because retrieval is frozen byte-for-byte across conditions, the ONLY thing that
%   varies between policies is what each chose to store and forget. Therefore any
%   between-policy difference in CL-F1 or footprint is attributable to storage alone —
%   this is the control that makes H1-H5 interpretable.
